%% ma: L11-B27-0002
%% lop: 11
%% chuong: 7
%% bai: 27
%% dang: 11.27.2
%% loai: DS
%% muc: [TH, TH, VD, TH]
%% dap_an: "ĐĐSĐ"
%% nguon: "(NBV)-ĐỀ SỐ 2-MỖI NGÀY 1 ĐỀ THI 2025.pdf (D02, MỖI NGÀY 1 ĐỀ - Đề số 2), Phần II Câu 1 (THPT Lương Tài số 2 - Bắc Ninh 2025)"
%% kien_thuc: "Bài 27 (thể tích khối chóp), Bài 26 (khoảng cách), Bài 24 (góc giữa đường thẳng và mặt phẳng), Bài 23 (đường thẳng vuông góc với mặt phẳng)"
%% ghi_chu: "Đề gốc ghi nhầm “Cho hình cho S.ABCD” (đã sửa thành hình chóp), không có hình trong đề. Đáp án gốc Đ Đ S Đ đúng (kiểm bằng sympy)"
%% trang_thai: chinh-thuc
\begin{ex}
Cho hình chóp $S.ABCD$ có đáy $ABCD$ là hình chữ nhật với $AB=2a$, $AD=6a$. Cạnh bên $SA$ vuông góc với mặt phẳng đáy, cạnh bên $SB=4a$.
\choiceTF
{\True $SA=2\sqrt3\,a$.}
{\True Góc giữa đường thẳng $SD$ và mặt phẳng $(ABCD)$ bằng $30^\circ$.}
{Khoảng cách từ điểm $C$ tới mặt phẳng $(SAD)$ bằng $6a$.}
{\True Thể tích của khối chóp $S.ABCD$ bằng $8\sqrt3\,a^3$.}
\loigiai{a) $SA\perp(ABCD)\Rightarrow SA\perp AB$, nên $SA=\sqrt{SB^2-AB^2}=\sqrt{16a^2-4a^2}=2\sqrt3\,a$. Đúng.
b) $AD$ là hình chiếu của $SD$ trên $(ABCD)$ nên góc cần tìm là $\widehat{SDA}$, $\tan\widehat{SDA}=\dfrac{SA}{AD}=\dfrac{2\sqrt3a}{6a}=\dfrac{\sqrt3}{3}\Rightarrow\widehat{SDA}=30^\circ$. Đúng.
c) $CD\perp AD$, $CD\perp SA$ nên $CD\perp(SAD)$, do đó $d(C,(SAD))=CD=AB=2a\ne6a$. Sai.
d) $V=\dfrac13\cdot AB\cdot AD\cdot SA=\dfrac13\cdot2a\cdot6a\cdot2\sqrt3a=8\sqrt3\,a^3$. Đúng.}
\end{ex}
